% =====================================================================
%  Präambel: Dokumentklasse, Pakete, Formatierung
% =====================================================================

\usepackage[utf8]{inputenc}
\usepackage[ngerman]{babel}
\usepackage{amsmath}
\usepackage{graphicx}
\usepackage{hyperref}
\usepackage{geometry}
\usepackage[absolute,overlay]{textpos}
\usepackage{setspace}

 %\usepackage{mathptmx}					% <--  Setzt Times New Roman als Schriftart

%	Provadis Richtlinien / Vorgaben:
%
%	- Seitenabstand: Oben: 3cm / Links: 3cm / Rechts: 2cm / Unten: 2cm
%
%	- Seitennummerierung (Seitenzahlen kommen immer in die Fußzeile): 
%		- Vor Textteil: Kleine röm. Zahlen (i,ii,iii,iv,v,...)
%		- Textteil: Arabische Zahlen (1,2,3,4,5,...)
%		- Nach Textteil: Große röm. Zahlen (I,II,III,IV,V,...)
%
%	- Schriftart (Größe):
%		- Überschriften: Times new Roman(14), Arial(13), Verdana(12)
%		- Text: Times new Roman(12), Arial(11), Verdana(10)
%		- Fußnoten: Times new Roman(10), Arial(9), Verdana(8)
%
%	- Deckblatt:
%		- Art d. Arbeit: bspw. Masterarbeit/Bachelorarbeit/Exposé
%		- Titel d. Arbeit
%		- Bei Praxisberichten, die Lehrveranstaltung, für die dies eine Prüfungsleistung ist
%		- Infos zum Autor: Name / Matrikelnummmer / Studiengang / Gruppe / E-mail
%		- Abgabedatum / Ende d. Bearbeitungsfrist
%		-Gutachter
%
%	- Verpflichtende Verzeichnisse:
%		- Inhaltsverzeichnis
%		- Literaturverzeichnis
%		- KI-Verzeichnis
%		- Ehrenwörtliche Erklärung
%
%	- Weitere Verzeichnisse (Nur bei bedarf verwenden):
%
%	- Rahmenbedingungen für d. Text:
%		- Text im Blocksatz formatieren
%		- Automatische Silbentrennzone: 1,25cm
%		- Zeilenabstand für Fließtext: 1,5 Zeilen
%
%
%
%

\geometry{top=3cm, bottom=2cm, left=3cm, right=2cm}

% Bilder werden im Verzeichnis assets/images/ gesucht
\graphicspath{{assets/images/}}

\title{Wissenschaftlich angeleitete Berufspraxis}
\date{}
